\documentclass[../../../include/open-logic-section]{subfiles}

\begin{document}

\olfileid{sth}{z}{pairs}
\olsection{जोड्या}

आता पुढील स्वयंसिद्धक पाहू:

\begin{axiom}[जोड्या]
कोणतेही संच $a, b$ दिल्यास $\{a, b\}$ हा संच अस्तित्वात असतो.
\[
	\forall a \forall b \exists P \forall x (x \in P \liff (x = a \lor x = b))
\]
\end{axiom}

क्रमिक संकल्पनेच्या आधारे या स्वयंसिद्धकाचे समर्थन असे करता येईल. $a$ टप्पा $S$ वर आणि $b$ टप्पा $T$ वर उपलब्ध आहेत असे मानू. $S$ आणि $T$ यांपैकी नंतर येणारा टप्पा $M$ असू द्या. मग $a$ आणि $b$ दोन्ही टप्पा $M$ वर उपलब्ध असल्याने, $\{a,b\}$ हा $M$ नंतरच्या कोणत्याही टप्प्यावर उपलब्ध वस्तूंपासून बनू शकणारा संग्रह आहे.

पण थांबा!{} $M$ नंतर टप्पे \emph{आहेतच} असे का मानावे? तसे नसतील तर आपले समर्थन अपुरे पडेल. म्हणून जोड्यांचे स्वयंसिद्धक समर्थित करण्यासाठी \olref[sth][z][story]{sec} मधील मांडणीत आणखी एक तत्त्व घालावे लागेल:
\begin{enumerate}
\item[] \stagessucc. शेवटचा टप्पा नसतो.
\end{enumerate}
या तत्त्वाचे समर्थन करता येते का? शोनफिल्ड यांच्या मांडणीत शेवटचा टप्पा नाही असे \emph{स्पष्टपणे} म्हटलेले नाही. तरी काटेकोरपणे पाहता हा अतिरिक्त भाग असला, तरी टप्प्याटप्प्याने संच तयार होतात या मूळ कल्पनेशी तो सुसंगत आहे. पुढील चर्चेत आपण हे तत्त्व स्वीकारू. त्याबरोबर जोड्यांचे स्वयंसिद्धकही स्वीकारू.

या नव्या स्वयंसिद्धकाच्या साहाय्याने आणखी अनेक संचांचे अस्तित्व सिद्ध करता येते. उदाहरणार्थ:

\begin{prop}\ollabel{prop:pairsconsequences}
कोणतेही संच $a$ आणि $b$ दिल्यास पुढील संच अस्तित्वात असतात:
\begin{enumerate}
\item\ollabel{singleton} $\{a\}$
\item\ollabel{binunion} $a \cup b$
\item\ollabel{tuples} $\tuple{a, b}$
\end{enumerate}
\end{prop}

\begin{proof}
\olref{singleton}. जोड्यांच्या स्वयंसिद्धकानुसार $\{a, a\}$ अस्तित्वात असतो; विस्तारात्मकतेनुसार तो $\{a\}$ च असतो.

\olref{binunion}. जोड्यांच्या स्वयंसिद्धकानुसार $\{a, b\}$ अस्तित्वात असतो. आता संयोगाच्या स्वयंसिद्धकानुसार $a \cup b = \bigcup
\{a, b\}$ अस्तित्वात असतो.

\olref{tuples}. \olref{singleton} नुसार $\{a\}$ अस्तित्वात असतो. जोड्यांच्या स्वयंसिद्धकानुसार $\{a,
b\}$ अस्तित्वात असतो. मग पुन्हा जोड्यांच्या स्वयंसिद्धकानुसार $\{\{a\}, \{a, b\}\} = \tuple{a, b}$ अस्तित्वात असतो.
\end{proof}

\begin{prob}
कोणतेही संच $a, b, c$ दिल्यास $\{a, b, c\}$ अस्तित्वात असतो, हे दाखवा.
\end{prob}

\begin{prob}
कोणतेही संच $a_1, \ldots, a_n$ दिल्यास $\{a_1, \ldots,
a_n\}$ अस्तित्वात असतो, हे दाखवा.
\end{prob}

%\begin{proof} By two applications of Pairs, $\{\{a_1, a_2\}, \{a_1,
%   a_3\}\}$ exists. By Union and Extensionality, $\bigcup \{\{a_1,
%   a_2\}, \{a_1, a_3\}\} = \{a_1, a_2, a_3\}$ exists. Repeat this
%   trick as often as necessary. \end{proof}

\end{document}
